\documentclass{article}
\usepackage[utf8]{inputenc}
\usepackage{amsmath}
\usepackage{geometry}
\geometry{margin=2cm}
\begin{document}

\section*{Análise da Questão 151 — ENEM 2024 — Matemática}

\subsection*{Enunciado}
O arquiteto Renzo Piano exibiu a maquete da nova sede do Museu Whitney de Arte Americana, um prédio assimétrico que tem um vão aberto para a galeria principal, cuja medida da área é 1.672 m$^{2}$. Considere que a escala da maquete exibida é 1 : 200. A medida da área do vão aberto nessa maquete, em centímetro quadrado, é?

\subsection*{Solução Técnica}

A resposta correta é a alternativa \textbf{E}, que corresponde à medida da área do vão aberto na maquete igual a 418,00 cm$^{2}$. A resolução envolve compreender a relação de escala entre a maquete e o prédio real, usando a escala $1:200$ para áreas.

\subsubsection*{Estratégia de Resolução}

\begin{enumerate}
  \item Identificar que a escala $1:200$ significa que 1 cm na maquete equivale a 200 cm no prédio real (escala linear).
  \item Converter a área real do vão (dada em m$^{2}$) para cm$^{2}$, sabendo que $1 \text{ m}^{2} = 10.000 \text{ cm}^{2}$.
  \item Determinar o fator de escala para a área, que é o quadrado do fator da escala linear:
  \[
    (\frac{1}{200})^{2} = \frac{1}{40.000}.
  \]
  \item Calcular a área na maquete multiplicando a área real em cm$^{2}$ pelo fator de escala para área.
\end{enumerate}

\subsubsection*{Cálculos}

Dado:

\begin{align*}
  \text{Área real} &= 1.672 \text{ m}^2, \\
  1 \text{ m}^2 &= 10.000 \text{ cm}^2 \implies \text{Área real em cm}^2 = 1.672 \times 10.000 = 16.720.000 \text{ cm}^2.
\end{align*}

Aplicando o fator da escala para áreas:

\[
  \text{Área na maquete} = 16.720.000 \times \frac{1}{40.000} = 418 \text{ cm}^2.
\]

Portanto, a medida da área do vão aberto na maquete é 418 cm$^{2}$.

\subsection*{Explicação Didática}

Para calcular a área do vão aberto na maquete, deve-se compreender que a escala \emph{linear} informa relação entre comprimentos, mas para áreas temos que elevar essa escala ao quadrado. Após converter a área real de metros quadrados para centímetros quadrados, aplicamos o fator da escala ao quadrado para obter a área correspondente na maquete.

\begin{enumerate}
  \item Entenda que escala $1:200$ significa que 1 cm na maquete equivale a 200 cm (2 metros) no prédio real.
  \item Converta a área real do vão $1.672$ m$^{2}$ para cm$^{2}$:
  \[
    1.672 \,\text{m}^{2} = 16.720.000 \,\text{cm}^{2}.
  \]
  \item Calcule a escala para área elevando ao quadrado:
  \[
    (1/200)^2 = 1/40.000.
  \]
  \item Multiplique a área real em cm$^{2}$ pela escala ao quadrado para encontrar a área na maquete:
  \[
    16.720.000 \,\text{cm}^{2} \times \frac{1}{40.000} = 418 \,\text{cm}^{2}.
  \]
  \item A área do vão aberto na maquete é 418 cm$^{2}$, que corresponde à alternativa \textbf{E}.
\end{enumerate}

\subsection*{Erros Comuns}

\begin{itemize}
  \item Usar a escala linear diretamente para a área, sem elevá-la ao quadrado.
  \item Não converter corretamente as unidades de medida antes de aplicar a escala.
  \item Confundir a escala linear com a escala para área.
  \item Erros de cálculo ao aplicar o fator da escala.
\end{itemize}

\subsection*{Dicas de Ensino}

\begin{itemize}
  \item Esclarecer a diferença entre escala linear e escala para área.
  \item Demonstrar a necessidade de converter unidades corretamente antes da aplicação da escala.
  \item Utilizar exemplos visuais para representar a relação entre escala linear e área.
  \item Incentivar reconhecimento do tipo de medida para aplicar a escala correta.
  \item Praticar exercícios envolvendo diferentes unidades e tipos de escala.
\end{itemize}

\subsection*{Análise dos Distratores}

\begin{description}
  \item[Alternativa A:] Usa a escala linear diretamente para a área (erro de escala dimensional).
  \item[Alternativa B:] Erro similar ao da alternativa A, aplicação incorreta da escala.
  \item[Alternativa C:] Erro em cálculo ou conversão parcial, valor não correto.
  \item[Alternativa D:] Confusão entre unidades e escala, resultado incorreto.
\end{description}

\subsection*{Perfil da Questão}

\begin{itemize}
  \item Habilidade principal: interpretação e aplicação da escala em figuras geométricas.
  \item Habilidades secundárias: conversão de unidades e cálculo com áreas.
  \item Demanda cognitiva: nível médio.
  \item Dificuldade estimada: média.
  \item Requer cálculo e conversão, não exige modelagem ou interpretação visual complexa.
  \item Observações: importância da compreensão da diferença entre escala linear e escalar, conversão adequada de unidades e aplicação correta da escala.
\end{itemize}

\subsection*{Revisão de Consistência}

A análise técnica está correta, as explicações estão coerentes e os distratores são adequados, apontando os principais erros comuns de entendimento de escala e unidades. Não apresenta incoerências relevantes.


\end{document}